In compliance with the course guidelines regarding academic integrity and content responsibility, we declare that we have used Generative Artificial Intelligence (GenAI) tools to support the drafting of this document (RASD).

We emphasize that the entire content of this document is the result of our original reasoning; AI was not used to generate architectural decisions or core logic from scratch, but exclusively as a tool for linguistic support, formal verification, and visual prototyping. The details of its usage are reported below:

\begin{enumerate}
    \item \textbf{Tools Used}
    \begin{itemize}
        \item \textbf{Textual Support:} ChatGPT (model GPT-3.5/4o mini), Gemini Advanced (Gemini 3).
        \item \textbf{Visual Generation (UI Mockups):} Lovable.
    \end{itemize}

    \item \textbf{Provided Inputs (Prompts and Context)} The inputs provided to the system were exclusively textual and based on preliminary work carried out by the group. Specifically:
    \begin{itemize}
        \item \textbf{For linguistic revision:} We input paragraphs or entire sections written by us in Italian/English, with the request to improve grammatical form, academic vocabulary, and text fluidity, without altering the technical meaning.
        \begin{quote}
            \textbf{Example A - For stylistic improvement:} \textit{"Act as a professional Technical Writer reviewing a Software Design Document. I will provide you with a raw text describing the 'Path Merging Algorithm'. Please rewrite it to improve grammatical correctness and use formal academic English suitable for a system architecture description. Do NOT add new information or change the logical flow. Just polish the language."}
        \end{quote}
        
        \item \textbf{For completeness verification:} We provided lists of component interfaces or sequence diagram steps developed by the group, asking the model to flag potential logical ambiguities or missing edge cases.
        \begin{quote}
            \textbf{Example B - For completeness verification:} \textit{"Act as a Senior Software Engineer. Below is the list of APIs we defined for the 'Sensor Controller'. Analyze them for consistency with the REST standard. Specifically, identify if there are any missing response codes or parameters that a standard Design Document should include."}
        \end{quote}
        
        \item \textbf{For UI Prototyping (via Lovable):} We provided detailed descriptions of the application screens, including layout structure, color palettes, and required UI elements, to generate high-fidelity visual mockups for Chapter 3.
        \begin{quote}
            \textbf{Example C - For UI Generation:} \textit{"Create a mobile screen for a cycling app in a modern style. It must include a map view showing a route, a bottom sheet showing real-time speed (25 km/h) and distance (12.4 km), and a large red circular button to 'Report Hazard'."}
        \end{quote}
    \end{itemize}

    \item \textbf{Obtained Outputs}
    \begin{itemize}
        \item Suggestions for sentence reformulation to make the text more formal and concise.
        \item Checklists to verify consistency between different diagrams and textual descriptions.
        \item Graphical assets and UI Mockups representing the application interface, which were subsequently refined to match the UX flow.
    \end{itemize}

    \item \textbf{Verification, Refinement, and Integration} Every output generated by the AI was subjected to a manual critical review by the group members.
    \begin{itemize}
        \item We did not passively accept the proposed reformulations; instead, we selected and adapted them to ensure they exactly reflected our project specifications.
        \item The UI mockups generated by Lovable were manually reviewed and used as a basis to define the final interface design, ensuring consistency with the functional requirements.
    \end{itemize}
\end{enumerate}

